\section{Implementation Details}
\label{app:implementation}

\begin{table*}[ht]
\centering
\caption{\textbf{Layered Memory System Construction via Standardized Formalization.} \plugmem adopts a layered memory graph: raw interactions ($\tau$) are first formalized into episodic memory (Phase 1). This formalized history then serves as the input for inducing semantic and procedural knowledge blocks (Phase 2). All memory types are managed through a universal API by standard data operations.}
\label{tab:memory_implementation}
\resizebox{\textwidth}{!}{%
\begin{tabular}{@{}l l l l@{}}
\toprule
\textbf{Memory Type} & \textbf{Construction Logic (Input $\rightarrow$ Output)} & \textbf{Graph Structure} & \textbf{Universal API \& Cognitive Heuristics} \\ \midrule

% --- Phase 1 Header ---
\multicolumn{3}{c}{\textbf{\textit{Phase 1: Standardize}}} & \multirow{14}{*}{\begin{tabular}[t]{@{}l@{}} 
\textbf{1. Create (Ingestion \& Induction)} \\ 
\quad $\bullet$ \texttt{memory.create(trajectory)} \\ 
\quad $\bullet$ Formalizes raw $\tau$ into standardized Episodic, and \\ 
\quad then triggers induction to abstract Facts/Skills from \\ 
\quad new Episodic entities.\\ 
\textbf{2. Retrieve (Association)} \\ 
\quad $\bullet$ \texttt{memory.retrieve(goal, state)} \\
\quad $\bullet$ Maps Goal and State to Intent/Concept indices \\ 
\quad and retrieve Knowledge Blocks, then traces back to \\ 
\quad Episodic Sources for context. \\ 
\textbf{3. Update (Reflection \& Refinement)} \\ 
\quad $\bullet$ \texttt{memory.update(feedback)} \\
\quad $\bullet$ Detect and mitigate contradictions (Semantic) or \\ 
\quad optimize workflow efficiency (Procedural). \\ 
\textbf{4. Delete (Forgetting \& Pruning)} \\ 
\quad $\bullet$ \texttt{memory.delete(criteria)} \\
\quad $\bullet$ Decays low-utility Nodes (Semantic/Procedural) \\ 
\quad  while compressing old Episodic Nodes to save space.\\ 
\end{tabular}} \\ \cmidrule(r){1-3}

% --- Episodic Row ---
\textbf{Episodic} & \begin{tabular}[c]{@{}l@{}}\textbf{Formalization} \\ \textit{Input:} Raw Agent Trajectory ($\tau$) \\ \textit{Process:} $Formalize(\tau) \rightarrow M_{epi}$ \\ \textit{Role:} Converts unstructured logs into \\ standardized Source Nodes.\end{tabular} & \begin{tabular}[c]{@{}l@{}}\textbf{Source Node (Event Window)} \\ \textit{Edge:} \texttt{Knowledge} $\xleftarrow{proves}$ \texttt{Source}\end{tabular} & \\ \cmidrule(r){1-3}

% --- Phase 2 Header ---
\multicolumn{3}{c}{\textbf{\textit{Phase 2: Structure}}} & \\ \cmidrule(r){1-3}

% --- Semantic Row ---
\textbf{Semantic} & \begin{tabular}[c]{@{}l@{}}\textbf{Fact Induction} \\ \textit{Input:} Episodic Memory ($M_{epi}$) \\ \textit{Process:} $Induce_{fact}(M_{epi}) \rightarrow M_{sem}$ \\ \textit{Role:} Distills static truths from \\ specific event traces.\end{tabular} & \begin{tabular}[c]{@{}l@{}}\textbf{Proposition Node (Fact Block)} \\ \textbf{Concept Node (Entity/Term)} \\\textit{Edge:} \texttt{Concept} $\xleftarrow{mentions}$ \texttt{Proposition}\end{tabular} & \\ \cmidrule(r){1-3}

% --- Procedural Row ---
\textbf{Procedural} & \begin{tabular}[c]{@{}l@{}}\textbf{Skill Induction} \\ \textit{Input:} Episodic Memory ($M_{epi}$) \\ \textit{Process:} $Induce_{skill}(M_{epi}) \rightarrow M_{proc}$ \\ \textit{Role:} Synthesizes reusable workflows \\ from successful history.\end{tabular} & \begin{tabular}[c]{@{}l@{}}\textbf{Prescription Node (Workflow Block)} \\ \textbf{Intent Node (User Goal/Task)}\\ \textit{Edge:} \texttt{Intent} $\xleftarrow{solves}$ \texttt{Prescription}\end{tabular} & \\ \bottomrule
\end{tabular}%
}
\end{table*}
\paragraph{Overview.} {This appendix provides implementation-oriented details that complement the high-level description in Section~\ref{subsec:structuring}. As summarized in Table~\ref{tab:memory_implementation}, the memory structuring module includes two main phases: Phase 1 ``Standardize'' and Phase 2 ``Structure.'' In Phase 1, to handle different types of episodic memory, \plugmem standardize the raw agent trajectory $\tau$ into uniformed episodic memory $M_{epi}$. Specifically, for each turns of observation and action $\{o_t, a_t\}$, we prompt LLM to derive the state $s_t$, the subgoal $g_t$ and the reward $r_t$. In Phase 2, \plugmem employs induction processes to distill this episodic memory $M_{epi}$ into Semantic Memory and Procedural Memory. These memory are the inserted into the memory graph, with a series of universal API that manage them.
%%%%%%%%%%%%%%%%%%%%%%%%%%%%%%%%%%%%%%%%%%%
% \newpage
\subsection{Episodic Standardization}
\label{app:implementation_standardize_episodic}

To standardize episodic memory at time $t$, we represent each turns of the raw agent trajectory as an observation-action pair $\{o_t,a_t\}$ and formalize it through a three-stage process:\\
We first derive the agent state $s_t$ using the previous state $s_{t-1}$, the previous action $a_{t-1}$, the current observation $o_t$, and the initial goal $G$. Incorporating $G$ ensures the state remains grounded in the agent's initial objective. The derivation is achieved by querying an LLM with the prompt in Listing~\ref{lst:apdx_prompt_get_state}. The resulting state serves as a concise context summary, effectively reducing overall context length for distilling procedural memory when pursuing long-context agentic tasks~\cite{wu2024longmemeval,zhou2024webarenarealisticwebenvironment}.\\
Next, we infer a subgoal $g_t$ from $s_t$, $o_t$, $a_t$ and $G$ to analyze how the agent decomposes the global task into incremental steps. Using the prompt in Listing~\ref{lst:apdx_prompt_get_subgoal}, this step serves two purposes: it enriches the information available for knowledge extraction and enables trajectory segmentation. Specifically, we can structure the raw trajectory into workflow by chunking episodic memory based on the cosine similarity of adjacent subgoal embeddings.\\
Finally, we determine the reward $r_t$ based on $s_t$, $a_t$, $g_t$ and $o_{t+1}$ using the prompt in Listing~\ref{lst:apdx_prompt_get_reward}. This reward quantifies the effectiveness of the action in achieving its intended goal, facilitating the extraction of quality-aware procedural memory.

\begin{lstlisting}[
  caption={Template for deriving the current state $s_t$ given previous state $s_{t-1}$, previous action $a_{t-1}$, and current observation $o_t$.},
  label={lst:apdx_prompt_get_state},
]
============================================================
Prompt Get_State
============================================================

You will receive four pieces of information:
Goal: The agent's current objective or task.
Previous State (at time t): A natural language summary describing the agent's context, history, and partial progress so far.
Action (at time t): The action the agent decided to take next, expressed in natural language.
Observation (at time t+1): The outcome or feedback resulting from that action.

Your task is to derive the new updated state---a coherent natural language summary that integrates all relevant information from the previous state, the action, and the new observation.

Steps to Follow:
Interpret the Inputs:
Examine the goal, the previous state, the action, and the observation to understand what has changed in the agent's situation and the detailed information about location and time.

Reason about the Update:
Describe the logical process by which the new state should differ from the previous one. Identify what progress has been made, what new information was gained, and how the context or focus may have shifted.

Generate the Updated State:
Write a clear and concise natural-language description summarizing the new state of the agent at time t+1. The new state should:
-- Include all the detailed information in the action and observation, especially information about location and time.
-- Integrate the outcome of the latest action and observation.

Output Format:
### Reasoning
(Explain step by step how the new state should be updated based on the inputs.)

### State
(Provide the final updated state summary here.)

---
Input:
Goal: {goal}
Previous State (at time t): {state}
Action (at time t): {action}
Observation (at time t+1): {observation}
\end{lstlisting}


\begin{lstlisting}[
  caption={Template for deriving the subgoal $g_t$ of turn $t$ given initiate goal $G$, current state $s_{t}$, observation $o_t$ and action $a_{t}$, as well as next observation $o_{t+1}$.},
  label={lst:apdx_prompt_get_subgoal},
]
============================================================
Prompt Get_Subgoal
============================================================

At time t, the agent takes an action based on its state, observation, and overall goal. Your task is to infer the subgoal-the immediate or intermediate objective-that best explains why the agent chose this action.

Use the following information as context: Overall Goal: {goal} Current State (summary of past context): {state} Current Observation: {observation} Action at time t: {action}

Step 1: Reasoning Analyze how the current state and observation relate to the overall goal. Explain how the given action helps the agent make progress toward that goal-possibly by achieving a smaller intermediate objective. Be explicit and causal: describe why this action makes sense given the context.

Step 2: Subgoal Inference After reasoning, infer the agent's likely subgoal-a short natural-language statement that describes the immediate purpose behind the action.

Output Format:

### Reasoning
[Your reasoning process-a few sentences explaining how the action relates to the goal and context]

### Subgoal
[A short sentence describing the inferred subgoal]
\end{lstlisting}



\begin{lstlisting}[
  caption={Prompt template for deriving return from procedural memories.},
  label={lst:apdx_prompt_get_return},
]

============================================================
Prompt Get_Return
============================================================

You will receive:\newline
Goal Description: A text explaining what the agent was trying to achieve.\newline
Process Description: A text describing the agent's actions, decisions, and progress toward that goal.\newline
\newline
Your task:\newline
Analyze the agent's process and determine how much of the goal was completed, considering the following aspects:\newline
\newline
Grading Criteria (Score 1--10):\newline
10: The agent fully accomplishes the goal with no significant omissions; actions are fully aligned with the goal.\newline
8--9: The agent completes most of the goal with only minor gaps; strong alignment but not perfect.\newline
6--7: Partial completion; the agent covers many key elements but leaves notable parts unfinished or poorly executed.\newline
4--5: Limited progress; the agent attempts the goal but completes less than half or does so in an ineffective way.\newline
2--3: Very little completion; actions barely connect to the goal or achieve only minimal results.\newline
1: No meaningful progress; actions do not contribute to achieving the goal at all.\newline
\newline
Important Instructions:\newline
Base the score only on completion level and alignment with the stated goal.\newline
Do not provide explanations or commentary unless requested.\newline
Output must follow the format below exactly.\newline
\newline
Output Format:\newline
\texttt{\#\#\# Score}\newline
[number from 1 to 10]\newline
\newline
Input:\newline
---\newline
Goal:\newline
\{subgoal\}\newline
Process:\newline
\{procedural\_memory\}

\end{lstlisting}

\begin{lstlisting}[
  caption={Template for deriving the reward $r_t$ of action $a_t$ given current state $s_{t}$, goal $g_t$ and action $a_{t}$, as well as next observation $o_{t+1}$.},
  label={lst:apdx_prompt_get_reward},
]
============================================================
Prompt Get_Reward
============================================================

You will be given:
Goal: the agent's overall objective.
State (at time t): what the agent knew and had done before taking the action.
Action (at time t): the single action the agent chose.
Observation (at time t + 1): the immediate outcome produced by that action.
Your task is to infer the reward - that is, a evaluation in natural language on how the agent's action contributes (positively or negatively) to achieving the overall goal, based on the resulting observation.
Follow these steps carefully:
1. Reasoning Process:
Explain how the action relates to the Goal given the State, and whether the Observation matches the expected helpful or unhelpful outcome.
Consider whether the action advances progress, causes setbacks, reveals new useful information, or wastes effort.
Summarize your reasoning about the causal contribution of the action to the goal.
2. Final Reward:
Use descriptive language to write a concise natural-language evaluation of the agent's action.
The reward should express how much and in what way the action helped or hindered achieving the goal.
Input:
Goal: {goal}
State (at time t): {state}
Action (at time t): {action}
Observation (at time t + 1): {observation}
Output format:
### Reasoning
[Your reasoning process here]
### Reward
[Natural-language reward statement that evaluate the agent's action]
some prompt
\end{lstlisting}




%%%%%%%%%%%%%%%%%%%%%%%%%%%%%%%%%%%%%%%%%%%
% \newpage
\subsection{Extract Knowledge}
\label{app:implementation_extract_knowledge}
\paragraph{Semantic}

Given a standardized episodic unit $e_t$, we extract semantic memories as proposition--tag pairs.
We query an LLM using the prompt in Listing~\ref{lst:apdx_prompt_get_semantic}, which returns up to $N_{\max}=10$ items in a fixed schema.
Each item contains a \texttt{Statement} field (a proposition in the main paper) and a \texttt{Tags} field (a list of associated concepts).
We deterministically parse the LLM output into $(p, \mathcal{T})$ pairs, where $p$ is the proposition text and $\mathcal{T}$ is the corresponding concept set.



We materialize these outputs into the semantic graph $G^{S}$ by creating a low-level proposition node for each $p$ and high-level concept nodes for each $c \in \mathcal{T}$, caching dense embeddings and metadata for all nodes.
We add two edge types: \textit{i)} \emph{membership} edges linking propositions to their concept tags ($p \rightarrow c$), and \textit{ii)} \emph{provenance} edges linking each proposition back to its source episodic unit in the episodic graph $G^{E}$ ($p \rightarrow e_t$).
When multiple propositions originate from the same episodic unit, we additionally connect them with \emph{sibling} edges to preserve local co-occurrence structure.


\begin{lstlisting}[
  caption={Prompt template for extracting semantic knowledge from episodic trajectories.},
  label={lst:apdx_prompt_get_semantic},
]
============================================================
Prompt Get_Semantic
============================================================

You are an expert at extracting precise, factual information from documents.
Your output must prioritize specificity, avoid ambiguity, eliminate redundancy, and strictly follow all formatting rules.

**CORE INSTRUCTIONS:**

1. Fact/Statement Extraction & Deduplication:
* Identify distinct, factual statements from the document.
* Resolving Vague References: If the subject of a statement is a pronoun or a vague description (e.g., 'the film', 'the band', 'the company', 'he', 'she', 'they'), you MUST rewrite the statement based on understanding of whole document so that the subject is a fully specified and concrete entity name taken from the document. You are NOT allowed to keep vague subjects in the final fact/statement.
* Example of Resolving Vague References:
BAD: 'The film was directed by xxx.', 'This movie is produced by xxx'
GOOD: 'Vaada Poda Nanbargal was directed by Manikai.', 'Vaada Poda Nanbargal is produced by xxx'
* Concrete Phrasing: Every statement MUST be phrased using explicit, identifiable names or titles. You are FORBIDDEN from using ANY vague references including but not limited to: 'the tour', 'the film', 'the movie', 'the band', 'it', 'he', 'she', or 'they'.
* Example of Concrete Phrasing:
BAD: 'The tour earned over $50 million.'
GOOD: 'NSYNC's Second II None Tour earned over $50 million.'
* Statement Length Policy (IMPORTANT): Each fact/statement does NOT have to be a single short sentence. A statement MAY be a compact multi-sentence block that groups tightly related information, but it must contain AT MOST 4 sentences total. These sentences should come from the original document material, potentially lightly edited ONLY to resolve vague references.
* Avoid Redundancy: You MUST merge similar or overlapping facts into single, comprehensive statements. Do NOT create multiple statements that repeat the same core information with minor variations.


2. Tag Generation:
* For each fact/statement, generate a list of tags.
* The number of tags per fact is flexible and should reflect the information density of the statement.
* Tags should cover key spans such as: entity names, years, numbers, nationalities, languages, genres, roles, object types, descriptive words, etc.
* The MAJORITY of tags SHOULD be SHORT TEXT SPANS copied VERBATIM from the statement (exact substrings). These tags are directly grounded in the surface form of the text.
* You MAY occasionally create additional tags that are not literal substrings of the current statement, in the following cases:
* You combine adjacent words into a single phrase (e.g., 'romantic comedy film').
* You import a surface span (exact phrase) from ANOTHER part of the document to make the meaning of the current statement explicit (for example, when resolving pronouns or phrases like 'the film', 'the band', etc.).


* When the subject of a statement is a pronoun or a vague description (e.g., 'the movie', 'the band', 'the hotel'), you MUST add at least one tag that names the underlying entity explicitly, using the exact surface form that appeared elsewhere in the document. This cross-statement tag may come from a previous or later sentence as long as it is clearly the same entity.
* If a tag is a verb, you MUST use its base (lemma) form (e.g., 'play', not 'played', 'playing', or 'to play').
* No Schema/Type Labels: You are FORBIDDEN from using meta-labels or ontology-like category names such as: 'Name', 'Person', 'Year', 'Date', 'City', 'Country', 'Location', 'Genre', 'Language', 'FilmTitle', etc. Tags are NOT type labels; they are content-bearing phrases.
* Tags should be relatively short and as fine-grained as needed. For example, adjectives and modifiers like 'Indian', 'Tamil-language', 'romantic', 'comedy' SHOULD usually be separate tags if they appear in the statement.



OUTPUT CONSTRAINTS:
* Extract up to 10 facts, but prioritize QUALITY over quantity. If there are fewer than 10 truly distinct facts, output fewer.
* Each fact must provide unique information not covered by other facts.
* ABSOLUTELY NO generic references -- every statement must explicitly name the specific entity.

**DOCUMENT:**
{observation}

OUTPUT FORMAT:

### Facts

1. **Statement:** [statement]
**Tags:** [tag0, tag1, tag2, tag3, ...]
2. **Statement:** [statement]
**Tags:** [tag0, tag1, tag2, tag3, ...]
...

\end{lstlisting}



\paragraph{Procedural} Given a standardized episodic sequence $M_{\mathrm{epi}}=[e_t]_{t=1}^{T}$ with subgoal annotations $g_t$, we first segment the full trajectory into coherent workflow chunks by detecting shifts in adjacent subgoals. We compute embeddings of each subgoal and start a new segment if the two adjacent subgoals have similarity below a set threshold. Each segment $M^{(i)}_{\mathrm{epi}}$ is then linearized into a compact process trace (e.g., a stepwise state--action--reward description) and fed to an LLM using the prompt in Table~\ref{lst:apdx_prompt_get_procedural} to produce a structured \emph{(intent, prescription)} pair $(u^{(i)}, \pi^{(i)})$, where $u^{(i)}$ abstracts the segment objective and $\pi^{(i)}$ describes an environment-agnostic procedure that can be reused in future tasks.

To support quality-aware reuse, we assign each prescription a scalar return score $\rho^{(i)}\in[1,10]$ via a separate LLM evaluator (Table~\ref{lst:apdx_prompt_get_return}), conditioned on the intent and the segment trace. We then materialize $(u^{(i)}, \pi^{(i)}, \rho^{(i)})$ into the procedural graph $G_P$ by creating an intent node and a prescription node, caching embeddings and metadata for both. To reduce intent proliferation, an incoming intent embedding $\phi(u)$ is matched against existing intent nodes; if the best cosine similarity exceeds threshold $\theta_{\mathrm{equal}}$, we merge the two intent strings via an LLM rewrite (Table~\ref{lst:apdx_prompt_get_new_subgoal}) and refresh the node embedding, otherwise we create a new intent node. The hierarchical relation is represented as a directed adjacency from intent to its prescriptions, i.e., $u \xrightarrow{\textsc{solves}} \pi$ (implemented as intent-node child lists, hence unidirectional). Finally, each prescription node is linked back to its originating episodic evidence through provenance edges $\pi \rightarrow e_t$ for all $e_t \in M^{(i)}_{\mathrm{epi}}$, enabling downstream verification and optional recovery of concrete interaction context when a retrieved procedure requires elaboration.
\begin{lstlisting}[
  caption={Prompt template for extracting procedural experience from episodic trajectories.},
  label={lst:apdx_prompt_get_procedural},
]
============================================================
Prompt Get_Procedural
============================================================

You will be given an utterance of an agent. Your task is to analyze the utterance and derive 
-- a main goal that the agent is pursuing. 
-- an experiential insight --- a concise reflection that summarizes the agent's behaviour.

Follow this process when generating your response:

Reasoning: Analyze the utterance and write down the generalizable information and patterns that would be useful as memory for future tasks.

Output goal and experiential insight: Produce one sentence describing the general goal of the utterance, using abstract language. Produce one paragraph in natural language that clearly expresses the experiential insight and the reflection that summarizes the agent's behaviour.

Output format:

### Reasoning
[Your reasoning process]

### Goal
[A sentence that concludes the goal]

### Experiential Insight
[The paragraph that expresses the experiential insight]

Input: Trajectory: {trajectory}


============================================================
Prompt Get_Return
============================================================

You will receive:
Goal Description: A text explaining what the agent was trying to achieve.
Process Description: A text describing the agent's actions, decisions, and progress toward that goal.

Your task:
Analyze the agent's process and determine how much of the goal was completed, considering the following aspects:

Grading Criteria (Score 1--10):
10: The agent fully accomplishes the goal with no significant omissions; actions are fully aligned with the goal.
8--9: The agent completes most of the goal with only minor gaps; strong alignment but not perfect.
6--7: Partial completion; the agent covers many key elements but leaves notable parts unfinished or poorly executed.
4--5: Limited progress; the agent attempts the goal but completes less than half or does so in an ineffective way.
2--3: Very little completion; actions barely connect to the goal or achieve only minimal results.
1: No meaningful progress; actions do not contribute to achieving the goal at all.

Important Instructions:
Base the score only on completion level and alignment with the stated goal.
Do not provide explanations or commentary unless requested.
Output must follow the format below exactly.

Output Format:
### Score
[number from 1 to 10]

Input:
---
Goal:
{subgoal}
Process:
{procedural_memory}
\end{lstlisting}


\begin{lstlisting}[
  caption={Prompt template for merging two similar intents.},
  label={lst:apdx_prompt_get_new_subgoal},
]
============================================================
Prompt Get_New_Subgoal
============================================================

Each goal may contain overlapping or complementary information.
Your task is to carefully combine them into a single, coherent, and well-structured goal that preserves all important details from both.
Avoid redundancy and ensure the merged goal sounds natural and consistent in tone.

Input:
Earlier goal: {goal_1}
Later goal: {goal_2}

Output:
Merged goal: [Write the unified goal here]
\end{lstlisting}





%%%%%%%%%%%%%%%%%%%%%%%%%%%%%%%%%%%%%%%%%%%
% \newpage
\subsection{Retrieval with Interleaved Abstraction and Specificity}
\label{app:implementation_retrieval}


This appendix provides implementation-oriented details that complement the high-level description in Section~\ref{subsec:retrieval}.
\plugmem performs interleaved multi-hop retrieval over the semantic graph $G^{S}$ and procedural graph $G^{P}$ by alternating between \textit{i)} abstraction-level routing through concept/intent nodes and \textit{ii)} specificity-level expansion over proposition/prescription nodes.
Given a query $Q$, the retriever is orchestrated by three LLM prompts: \textsc{GetMode} (Listing~\ref{lst:apdx_prompt_get_mode}) selects the memory type(s) to emphasize; \textsc{GetPlan} (Table~\ref{lst:apdx_prompt_get_plan}) proposes the next-hop abstract signals (concept tags for $G^{S}$ or subgoals/intents for $G^{P}$); and \textsc{MultiHopCtrl} (Listing~\ref{lst:apdx_prompt_multi_hop_retrieval}) serves as a retrieval controller that decides whether the current evidence is sufficient to stop, or which low-level nodes to prioritize in the next hop.



At hop $t$, retrieval starts from a low-level candidate pool $C_t$.
We first expand candidates through two complementary channels.
\emph{(Embedding channel)} We embed the current query text $Q_t$ and retrieve a set of high-scoring low-level nodes by dense similarity.
\emph{(Linking channel)} In parallel, we run \textsc{GetPlan} to produce an abstract plan (concept tags for $G^{S}$ or intents/subgoals for $G^{P}$), match it to the corresponding high-level nodes, and activate their adjacent low-level neighbors via membership edges.
The two candidate sets are merged (union), then reranked and pruned back to a fixed budget (top-$K$) using a value function (primarily relevance, optionally combined with metadata such as importance/recency).
Throughout the process, high-level nodes are used only as intermediate routing signals and are not retained in $C_t$.



After obtaining the budgeted candidate set, we invoke \textsc{MultiHopCtrl} to assess sufficiency.
Concretely, \textsc{MultiHopCtrl} returns a strict JSON object with an \texttt{enough} flag and an optional list of selected low-level node IDs (\texttt{top\_node\_ids}), constrained to be a subset of the available candidates; if \texttt{enough=true}, it returns an empty list.
If \texttt{enough=true}, retrieval terminates early and the current candidates are returned (or mapped to provenance-linked episodic nodes when episodic memory is requested).
Otherwise, we keep only the selected nodes (capped in our experiments) as the ``focus set'' and integrate them with the previous query to form the next-hop query $Q_{t+1}$.
This integration can be performed by an LLM-based rewriter, but we find that simple concatenation of $Q_t$ and the selected facts performs comparably well and is used as the default for efficiency.


The retrieval loop repeats until \textsc{MultiHopCtrl} signals sufficiency or a hop limit is reached.
The overall control flow is summarized in Listing~\ref{lst:apdx_multi_hop_retrieval_pseudo_code}.


\begin{lstlisting}[
  caption={Prompt template for inferring potential memory type to be used from query.},
  label={lst:apdx_prompt_get_mode},
]
============================================================
Prompt Get_Mode
============================================================

You are given a task description that the agent is pursuing and the observation from the task
Please analyze the task description and observation to determine the type of memory required to complete it effectively. There are three possible memory types:
Episodic Memory: This is needed if the task requires you to answer questions based on events. For example, answering user's question depending on historical conversation.
Semantic Memory: This is needed if the task requires you to recall objective information. For example, answer the question based on objective knowledge or information.
Procedural Memory: This is needed if the task is completing a subgoal under an interactive environment that agent need to perform a workflow. For example, completing an instruction in web navigation tasks.
First analyze that task and observation and decide which only one memory type needed. When there is a conflict, prioritize the information in the Task Description when making decisions.

Output Format:

### Reasoning

[Your analyze of which memory is needed depending on task and observation]

### Memory Type

## [Your final decision, episodic_memory or semantic_memory or procedural_memory]

Input:
Task Description: {task_type}
Observation: {observation}
\end{lstlisting}

\begin{lstlisting}[
  caption={Prompt template for inferring abstraction concepts for next step of retrieval.},
  label={lst:apdx_prompt_get_plan},
]
============================================================
Prompt Get_Plan
============================================================

You are an expert at analyzing an agent's goal and current observation and generating retrieval tags for a goal-directed question-answering system.
Your setting:

* Goal: The agent's overall objective to accomplish.
* Current Subgoal: The subgoal the agent is currently pursuing. (can be None)
* Current State: A description of the agent's current internal state. (can be None)
* Input (Current Observation): The agent's most recent observation (typically a question or task instruction).
* Task: Extract a prioritized set of high-quality tags that are most likely to retrieve information that directly helps accomplish the Goal.

Instructions:

1. Goal-directed Tag Selection (CRITICAL):

* Read the Goal and the Current Observation carefully.
* Only generate tags that are HIGHLY LIKELY to retrieve evidence needed to solve/complete the Goal.
* Prefer tags that identify:
* The target entity/entities the Goal is asking about (people, organizations, places, works, events).
* Bridge entities implied by the observation that are likely required for multi-hop retrieval.
* Explicit constraints: dates, years, roles, titles, unique descriptors, numbers.


* Avoid low-signal or generic tags that are unlikely to retrieve helpful evidence (e.g., "known for", "famous", "character" unless the Goal specifically depends on them).

2. Concrete, Grounded Tags:

* The MAJORITY of tags MUST be short text spans copied VERBATIM from the Goal or the Current Observation (exact substrings).
* You MAY add a SMALL number of non-literal tags only if they are short, strongly implied, and clearly necessary for retrieval (e.g., a canonical name expansion or a standard alias).
* If a tag is a verb, you MUST use its base (lemma) form (e.g., "direct", not "directed" or "directing").

3. Prioritization & Quantity:

* the number of tags in total should be proper. (not "as many as possible").
* Sort tags by expected retrieval usefulness (most useful first).
* Ensure tags are content-bearing and relatively short.

4. CRITICAL -- Forbidden tags:

* Do NOT generate the tag "user".
* Do NOT use meta-labels or type names such as "Name", "Person", "Year", "Date", "City", "Country", "Location", "Genre", "FilmTitle", etc.

Output Format:

### Reasoning

[You process of analyzing the information and completing the task]

### Tags

**Tags:** ["tag0", "tag1", "tag2", "tag3", ...]
(for example: "Central Area", "focal point", "famous for", etc)

### Next Subgoal

## [A single best next subgoal that the agent should pursue now.]

Input:
Goal: {goal}
Current Subgoal: {subgoal}
Current State: {state}
Current Observation: {observation}
\end{lstlisting}


\begin{lstlisting}[
  caption={Prompt template for controlling mutli-hop retrieval process.},
  label={lst:apdx_prompt_multi_hop_retrieval},
]
============================================================
Prompt Multi-hop_Retrieval
============================================================

You are a retrieval controller for multi-hop question answering.
Return **STRICT JSON only**:
{
"enough": true/false,
"top_node_ids": [int, int, ...]
}

**Constraints:**
-- top_node_ids length <= {n_facts_new_query}
-- top_node_ids must be a subset of available ids
-- if enough=true => top_node_ids=[]

**Question:**
{question}

**Available node ids:**
{available_ids}

**Retrieved facts:**
{semantic_memory_str}

\end{lstlisting}


% \begin{lstlisting}[
%   caption={Multi-hop retrieval algorithm procedure.},
%   label={lst:apdx_multi_hop_retrieval_pseudo_code},
% ]

% \ttfamily
% \textbf{Input:} query $Q$; graphs $(G^{E},G^{S},G^{P})$ with provenance links to $G^{E}$; budget $K$; hop limit $T_{\max}$\\
% \textbf{Output:} retrieved context $\mathcal{R}$\\[2pt]

% $\mathcal{F} \leftarrow \mathrm{TaskAdapter}(Q)$ \hfill \{ $\mathcal{F}\subseteq\{E,S,P\}$ \}\\
% $C \leftarrow \mathrm{InitLowLevel}(Q,\{G^{S},G^{P}\})$ \hfill \{ score all low-level nodes \}\\
% $C \leftarrow \mathrm{TopK}(C,K)$\\[2pt]

% \textbf{for} $t=1$ \textbf{to} $T_{\max}$ \textbf{do}\\
% \hspace*{1.5em}\textbf{if} $\mathrm{Stop}(Q,C)$ \textbf{then break}\\
% \hspace*{1.5em}$q^{a} \leftarrow \mathrm{Abstract}(Q,C)$ \hfill \{ concepts for $G^{S}$; intents for $G^{P}$ \}\\
% \hspace*{1.5em}$C \leftarrow C \cup \mathrm{ExpandLowLevel}(q^{a},\{G^{S},G^{P}\})$\\
% \hspace*{1.5em}$C \leftarrow \mathrm{RerankPrune}(C,K)$ \hfill \{ e.g., relevance/importance/recency \}\\
% \textbf{end for}\\[2pt]

% \textbf{if} $E \in \mathcal{F}$ \textbf{then}\\
% \hspace*{1.5em}$\mathcal{R} \leftarrow \mathrm{ToEpisodic}(C,G^{E})$ \hfill \{ map via provenance \}\\
% \textbf{else}\\
% \hspace*{1.5em}$\mathcal{R} \leftarrow C$\\
% \textbf{end if}\\
% \textbf{return} $\mathcal{R}$\\
% \normalfont

% \end{lstlisting}

\begin{table}[t]
\centering
\caption{Multi-hop retrieval control flow (pseudo-code style).}
\label{lst:apdx_multi_hop_retrieval_pseudo_code}

\begin{tcolorbox}[colback=white,colframe=black!30,boxrule=0.6pt,left=6pt,right=6pt,top=4pt,bottom=4pt,
width=0.8\linewidth]
% ]
\ttfamily
\textbf{Input:} query $Q$; graphs $(G^{E},G^{S},G^{P})$ with provenance links to $G^{E}$; budget $K$; hop limit $T_{\max}$\\
\textbf{Output:} retrieved context $\mathcal{R}$\\[2pt]

$\mathcal{F} \leftarrow \mathrm{TaskAdapter}(Q)$ \hfill \{ $\mathcal{F}\subseteq\{E,S,P\}$ \}\\
$C \leftarrow \mathrm{InitLowLevel}(Q,\{G^{S},G^{P}\})$ \hfill \{ score all low-level nodes \}\\
$C \leftarrow \mathrm{TopK}(C,K)$\\[2pt]

\textbf{for} $t=1$ \textbf{to} $T_{\max}$ \textbf{do}\\
\hspace*{1.5em}\textbf{if} $\mathrm{Stop}(Q,C)$ \textbf{then break}\\
\hspace*{1.5em}$q^{a} \leftarrow \mathrm{Abstract}(Q,C)$ \hfill \{ concepts for $G^{S}$; intents for $G^{P}$ \}\\
\hspace*{1.5em}$C \leftarrow C \cup \mathrm{ExpandLowLevel}(q^{a},\{G^{S},G^{P}\})$\\
\hspace*{1.5em}$C \leftarrow \mathrm{RerankPrune}(C,K)$ \hfill \{ e.g., relevance, etc. \}\\
\textbf{end for}\\[2pt]

\textbf{if} $E \in \mathcal{F}$ \textbf{then}\\
\hspace*{1.5em}$\mathcal{R} \leftarrow \mathrm{ToEpisodic}(C,G^{E})$ \hfill \{ map via provenance \}\\
\textbf{else}\\
\hspace*{1.5em}$\mathcal{R} \leftarrow C$\\
\textbf{end if}\\
\textbf{return} $\mathcal{R}$
\end{tcolorbox}

\end{table}



%%%%%%%%%%%%%%%%%%%%%%%%%%%%%%%%%%%%%%%%%%%
% \newpage
\subsection{Reasoning}

We design tailored reasoning prompts for each memory type based on their unique structural characteristics.\\
In alignment with the objective of the reasoning module is to aggregate and condense retrieved memory mentioned in Section~\ref{subsec:reasoning}, we query the LLM to extract information relevant to the question, utilizing the prompt structure detailed in Listing~\ref{lst:apdx_prompt_reasoning_semantic}.\\
Given that episodic memory is often long and unstructured, we observed that information extraction is inefficient as most of the LLM is not deliberately trained on the task. Instead, we find that, like prompt in Listing~\ref{lst:apdx_prompt_reasoning_episodic}, directly asks the LLM to reason over every memory and answer the question successfully leveraging the reasoning ability of LLM while containing most of the useful information in the reasoning process. The entire output including the reasoning process is then preserved as the extracted information.\\
Procedural memory is primarily utilized during multi-turn decision-making tasks. To support that, we prompt the LLM to integrate diverse experiences consisting of both successful and failed memory from similar tasks. As shown in Listing~\ref{lst:apdx_prompt_reasoning_procedural}, this process generates coherent and actionable guidance to fuel the agent's next action.

\label{app:implementation_reasoning}
\begin{lstlisting}[
  caption={Prompt template for reasoning and compressing for retrieved episodic memories.},
  label={lst:apdx_prompt_reasoning_episodic},
]
============================================================
Prompt Reason_Episodic
============================================================

I will give you information of history chats between you and a user. Please answer the question based on the information. Answer the question step by step: first extract all the relevant information, and then reason over the memory to get the answer.

Information:

{information}

Current Date: {time}
Question: {question}
Answer (step by step):
\end{lstlisting}


\begin{lstlisting}[
  caption={Prompt template for reasoning and compressing for retrieved semantic memories.},
  label={lst:apdx_prompt_reasoning_semantic},
]
============================================================
Prompt Reason_Semantic
============================================================

I will give you several retrieved facts. Extract all the useful information relevant to the question.
In the output reasoning information, use the original wording from the retrieved facts as much as possible, and do not replace it with synonyms or near-synonyms.
If no useful information found, just return "null".

## Output format:

### Reasoning

(process of extract information)

### Information

## (The useful information you extract)

Input:
Facts: {semantic_memory}
Current Date: {time}
Question: {observation}
\end{lstlisting}

\begin{lstlisting}[
  caption={Prompt template for reasoning and compressing for retrieved procedural memories.},
  label={lst:apdx_prompt_reasoning_procedural},
]
============================================================
Prompt Reason_Procedural
============================================================

Following information will be provided:
Question: The question the user is asking.
Information: Several pieces of information that may be relevant to the question.

Your task:

1. Carefully read the user's question.
2. Analyze each piece of retrieved information and determine how relevant and useful it is for answering the question.
3. Based on your analysis, integrate all the useful information into a single coherent piece of content that helps the agent to answer the question. When you think information is insufficient or contradictory, generate the most possible information. The integrated content should be concise, accurate, and relevant to the user's question.

## Output format:

### Reasoning

(Your reasoning for analysis of given question and information.)

### Final Information

## (The synthesized information that should be provided to the agent.)

Input:
Question: {observation}
Information: {procedural_memory}
\end{lstlisting}





%%%%%%%%%%%%%%%%%%%%%%%%%%%%%%%%%%%%%%%%%%%
% \newpage
\subsection{Full Operations on Memory Graph}

\label{app:full_operation}


\paragraph{Setup.} We conduct multiple rounds of controlled experiments on the HotpotQA-induced semantic subgraph. We choose HotpotQA for two practical reasons: \textit{i)} its evidence is largely factual and extractive, which makes semantic memories comparatively stable and less sensitive to stylistic variance; and \textit{ii)} the resulting subgraph is lightweight enough to support fast iteration over different merge thresholds and update policies.

Given a semantic node, our update routine first performs \emph{candidate discovery} by collecting all semantic nodes that share at least one concept/tag with it (i.e., neighbors under the tag-induced projection). We then compute embedding similarity between the current node and each candidate, rank candidates by similarity, and select the top-$m$ candidates whose similarity exceeds a pre-defined threshold $\tau$ (we use $m{=}1$ by default, i.e., only the most similar candidate above threshold is considered). If a pair $(s_i, s_j)$ is selected, we trigger a merge operation by calling an LLM to \textit{i)} synthesize a new semantic memory that better summarizes the combined information, and \textit{ii)} decide whether to deactivate the original nodes (delete neither, delete both, etc.) according to explicit rules encoded in the prompt; the full prompt is shown in Listing~\ref{lst:apdx_prompt_merge_semantic}.

\begin{lstlisting}[
  caption={Prompt template for checking and merging semantic nodes in memory update.},
  label={lst:apdx_prompt_merge_semantic},
]
============================================================
Prompt Merge_Semantic
============================================================

You are given two memory items about a related topic. One came earlier (Information 1) and the other came later (Information 2).

Your tasks:
(1) Merge them into ONE improved, clear, concise statement. Do not invent new facts.
(2) Decide whether to deactivate (soft delete) the original two nodes after merging.

Inputs:
Information 1 (Earlier Information): {memory_earlier}
Information 2 (Later Information): {memory_later}

Deactivation decision rules (choose exactly ONE case):

Case A: "UPDATE_SAME_FACT"

* Condition: Information 1 and 2 are essentially describing the same fact/event, and Information 2 mainly updates/corrects/refines details of Information 1.
* Action: deactivate BOTH originals (earlier and later) because the merged node fully supersedes them.

Case B: "SAME_TOPIC_MERGE_WELL"

* Condition: Information 1 and 2 are strongly related under the same topic, and the merged statement reads naturally as a unified summary (not an awkward splice).
* Action: deactivate BOTH originals (earlier and later).

Case C: "WEAK_RELATED_STITCH_RISK"

* Condition: Information 1 and 2 are only weakly related; merging feels like stitching two segments; and deactivating either original would likely harm future retrieval.
* Action: deactivate NEITHER original.

Hard constraints:

* Output MUST be valid JSON (no Markdown, no extra text).
* relationship MUST be one of the three labels above.
* If relationship is Case A or B => deactivate_earlier=true AND deactivate_later=true.
* If relationship is Case C => deactivate_earlier=false AND deactivate_later=false.
* If the two memories conflict, prefer Information 2 as the more up-to-date.
* Output the simple reasoning of why you made the decision.

Output MUST be valid JSON with exactly these keys:

* merged_statement (string)
* relationship ("UPDATE_SAME_FACT" | "SAME_TOPIC_MERGE_WELL" | "WEAK_RELATED_STITCH_RISK")
* deactivate_earlier (boolean)
* deactivate_later (boolean)
* simple_reasoning (string)

Return ONLY the JSON object.
\end{lstlisting}





\paragraph{Thresholds and evaluation.}
We evaluate two merge thresholds, $\tau \in \{0.6, 0.7\}$, on the same HotpotQA subset. We report downstream QA performance (EM/F1) as well as graph statistics before and after update.
The results show that performance remains within normal fluctuation, while the graph becomes more compact with more controllable candidate fan-out. Concretely, with $\tau{=}0.6$, $477/3413$ semantic nodes trigger a merge; with $\tau{=}0.7$, $171/3413$ nodes trigger a merge, reflecting the expected trade-off between aggressiveness and conservativeness.

We compare the downstream QA performance on the same HotpotQA subset before and after semantic graph update under two merge thresholds. Without update, the system achieves an EM/F1 of $61.00/74.39$. With update at $\tau{=}0.6$, the EM/F1 becomes $63.00/73.97$, and with update at $\tau{=}0.7$, the EM/F1 becomes $62.00/74.65$. Overall, these differences are within normal run-to-run fluctuation, suggesting that semantic graph update/merge does not materially degrade end-to-end performance on this subset.

\paragraph{Graph quality: compactness and candidate controllability.}
We characterize the structural effect of update/merge along two axes.
\begin{itemize}
    \item \textbf{Compactness.}
We measure compactness using the number of active semantic nodes $N_s$, the number of used tags $N_t$, and the number of semantic--tag bipartite edges $E_{\text{bip}}$ (attachments).
With the stricter threshold $\tau{=}0.7$, the updated graph reduces active semantic nodes from $3413$ to $3242$ ($-5.0\%$), used tags from $12501$ to $11812$ ($-5.5\%$), and bipartite edges from $23230$ to $20604$ ($-11.3\%$), indicating lower redundancy in semantic memories and their tag attachments.
We also report the (non-deduplicated) upper bound of tag-induced semantic co-occurrence pairs, $\sum_{t} \binom{\deg(t)}{2}$, which decreases from $78279$ to $57581$ ($-26.5\%$). This suggests that update/merge effectively prunes redundant co-attachments that would otherwise create dense but uninformative connectivity.
  \item \textbf{Candidate controllability.}
Our retrieval/update pipeline uses tag-induced candidate generation: for a semantic node, we take the union of semantic IDs attached to its tags as the candidate set. To quantify how controllable the candidate fan-out is, we sample active semantic nodes and compute the size of their candidate sets (unique neighbors reachable via tags, excluding itself).
With $\tau{=}0.7$, the candidate set size becomes substantially smaller and thus more controllable: the sampled mean decreases from $38.36$ to $31.28$ ($-18.5\%$), and the sampled median decreases from $18.0$ to $14.0$ ($-22.2\%$). Intuitively, by merging near-duplicate semantic items and deactivating redundant originals, the graph reduces unnecessary expansions through shared tags, improving computational stability without degrading downstream EM/F1.
  
\end{itemize}


Overall, these results indicate that semantic update/merge can improve graph compactness and reduce redundancy while preserving retrieval support, as evidenced by stable HotpotQA accuracy under both thresholds.